\section*{Capítulo \ref{ch:g4:symmetry} --- La simetría}

\begin{solution}{exo:g4:symmetry:1}
Cuadrado: sí ($4$ ejes). Rectángulo: sí ($2$). Triángulo escaleno: no hay
doblez que funcione. Círculo: sí --- cualquier recta que pase por el centro.
\end{solution}

\begin{solution}{exo:g4:symmetry:2}
Cuadrado: $4$ (dos rectas medias y dos diagonales). Rectángulo: $2$ (solo las
rectas medias). Triángulo equilátero: $3$ (una por cada
\omterm{def:g4:geometry:polygon}{vértice}).
\end{solution}

\begin{solution}{exo:g4:symmetry:3}
Al doblar por la diagonal, las dos \omterm{def:g3:division:half}{mitades} triangulares tienen la misma
forma, pero no se superponen --- una está «girada del
revés»: el calco cae al lado de la figura, no encima. Así que la
diagonal no es un eje.
\end{solution}

\begin{solution}{exo:g4:symmetry:4}
La L imagen queda a $2$ casillas al \emph{otro} lado del eje,
escrita del revés (una L reflejada), con los mismos tamaños.
\end{solution}

\begin{solution}{exo:g4:symmetry:5}
La figura completa muestra la bandera y su reflejo boca abajo
por debajo del eje, mástil bajo mástil.
\end{solution}

\begin{solution}{exo:g4:symmetry:6}
$P$ no se mueve: al doblar por el eje, todo punto que está
\emph{sobre} el eje se queda exactamente donde está --- $P$ es su propia imagen.
\end{solution}

\begin{solution}{exo:g4:symmetry:7}
Eje vertical: A, H, O, T. Eje horizontal: B, E, H, O. Los dos: H, O.
Ninguno: F, N.
\end{solution}

\begin{solution}{exo:g4:symmetry:8}
Sí: el triángulo es isósceles y su eje pasa por el \omterm{def:g4:geometry:polygon}{vértice}
que hay entre los dos lados de $5$ cm y por el punto medio de la base de $3$ cm. El
doblez intercambia los dos lados iguales.
\end{solution}

\begin{solution}{exo:g4:symmetry:9}
Vale cualquier figura con exactamente dos ejes: un rectángulo (no cuadrado), una
cruz con brazos de dos longitudes distintas (horizontal larga y vertical
corta) o un óvalo. Los dos ejes se cortan siempre en el centro de la
figura.
\end{solution}

\begin{solution}{exo:g4:symmetry:10}
La casa completa es la \omterm{def:g3:division:half}{mitad} izquierda más su gemela reflejada: una pared
del doble de ancho que la media pared, bajo un tejado triangular entero. Cada
longitud de la derecha es igual a su pareja de la izquierda --- la simetría
copia las longitudes exactamente.
\end{solution}

\begin{solution}{exo:g4:symmetry:11}
Los dibujos respaldan las dos direcciones: el triángulo isósceles se dobla
por la recta que pasa por su \omterm{def:g4:geometry:polygon}{vértice} (dos lados iguales $\to$ eje); el
equilátero se dobla de tres maneras; el escaleno no tiene ningún eje (sin
lados iguales $\to$ sin eje). La doble afirmación de Aya es correcta --- se
demostrará en el \cref{ch:g6:symmetry}.
\end{solution}
